
\begin{explanation}
\textbf{Understanding Logistic Regression MLEs:}

The MLEs arise from the invariance property: since $\hat{\pi}_0 = a/(a+b)$ and $\hat{\pi}_1 = c/(c+d)$ maximize the likelihood, and logit is a one-to-one transformation, the MLEs for $\beta_0$ and $\beta_1$ follow directly. The odds ratio $OR = \exp(\beta_1) = \frac{bc}{ad}$ measures association: $OR=1$ means no association, $OR>1$ positive association, $OR<1$ negative association.
\end{explanation}

\subsubsection{Part (b): Variance and Confidence Interval}

\begin{solution}
We derive $\text{Var}(\hat{\beta}_1)$ using the delta method.

\textbf{Delta Method Application:}

Since $\hat{\beta}_1 = \log b + \log c - \log a - \log d$, the partial derivatives are:
\[
\frac{\partial \hat{\beta}_1}{\partial a} = -\frac{1}{a}, \quad \frac{\partial \hat{\beta}_1}{\partial b} = \frac{1}{b}, \quad \frac{\partial \hat{\beta}_1}{\partial c} = \frac{1}{c}, \quad \frac{\partial \hat{\beta}_1}{\partial d} = -\frac{1}{d}
\]

For large samples with $\text{Var}(a) \approx a$, $\text{Var}(b) \approx b$, $\text{Var}(c) \approx c$, $\text{Var}(d) \approx d$:

\begin{align*}
\text{Var}(\hat{\beta}_1) &\approx \left(-\frac{1}{a}\right)^2 a + \left(\frac{1}{b}\right)^2 b + \left(\frac{1}{c}\right)^2 c + \left(-\frac{1}{d}\right)^2 d \\
&= \frac{1}{a} + \frac{1}{b} + \frac{1}{c} + \frac{1}{d}
\end{align*}

\[
\boxed{\text{Var}(\hat{\beta}_1) \approx \frac{1}{a} + \frac{1}{b} + \frac{1}{c} + \frac{1}{d}}
\]

\textbf{95\% Wald CI for $\beta_1$:}

With $SE(\hat{\beta}_1) = \sqrt{\frac{1}{a} + \frac{1}{b} + \frac{1}{c} + \frac{1}{d}}$ and $z_{0.975} = 1.96$:

\[
\boxed{CI_{95\%}(\beta_1) = \hat{\beta}_1 \pm 1.96 \cdot SE(\hat{\beta}_1)}
\]

\textbf{95\% Wald CI for Odds Ratio $\exp(\beta_1)$:}

By exponentiating the CI for $\beta_1$:

\[
\boxed{CI_{95\%}(OR) = \left[\exp\left(\hat{\beta}_1 - 1.96 \cdot SE\right), \exp\left(\hat{\beta}_1 + 1.96 \cdot SE\right)\right]}
\]
\end{solution}

\begin{explanation}
\textbf{Delta Method Intuition:}

The variance formula shows smaller cell counts increase uncertainty. Zero cells make variance infinite (why all cells must be $>0$). Each cell contributes $1/\text{count}$ to the total variance. The Wald CI assumes normality of $\hat{\beta}_1$ for large samples.
\end{explanation}

\subsubsection{Part (c): Numerical Example}

\begin{solution}
Given: $a = 18$, $b = 12$, $c = 30$, $d = 10$.

\textbf{Contingency Table:}

\begin{center}
\begin{tabular}{c|cc|c}
& $Y=1$ & $Y=0$ & Total \\
\hline
$X=0$ & 18 & 12 & 30 \\
$X=1$ & 30 & 10 & 40 \\
\hline
Total & 48 & 22 & 70
\end{tabular}
\end{center}

\textbf{Step 1: Compute $\hat{\beta}_0$}

\begin{align*}
\hat{\beta}_0 &= \log\frac{18}{12} = \log 1.5 \approx 0.4055
\end{align*}

\[
\boxed{\hat{\beta}_0 \approx 0.4055}
\]

The odds of $Y=1$ for $X=0$ are $\exp(0.4055) = 1.5$ (i.e., 3:2 odds).

\textbf{Step 2: Compute $\hat{\beta}_1$}

\begin{align*}
\hat{\beta}_1 &= \log\frac{12 \times 30}{18 \times 10} = \log\frac{360}{180} = \log 2 \approx 0.6931
\end{align*}

Verification: $\log(30/10) - \log(18/12) = \log 3 - \log 1.5 = \log 2$ ✓

\[
\boxed{\hat{\beta}_1 = \log 2 \approx 0.6931}
\]

\textbf{Step 3: Odds Ratio}

\[
\widehat{OR} = \exp(0.6931) = 2.0
\]

\[
\boxed{\widehat{OR} = 2.0}
\]

The odds of $Y=1$ for $X=1$ are twice those for $X=0$.

\textbf{Step 4: Variance and SE}

\begin{align*}
\text{Var}(\hat{\beta}_1) &= \frac{1}{18} + \frac{1}{12} + \frac{1}{30} + \frac{1}{10} \\
&= 0.0556 + 0.0833 + 0.0333 + 0.1000 = 0.2722
\end{align*}

\[
SE(\hat{\beta}_1) = \sqrt{0.2722} \approx 0.5218
\]

\[
\boxed{SE(\hat{\beta}_1) \approx 0.522}
\]

\textbf{Step 5: 95\% CI for $\beta_1$}

Margin of error: $1.96 \times 0.5218 = 1.023$

\begin{align*}
CI_{95\%}(\beta_1) &= [0.693 - 1.023, \, 0.693 + 1.023] \\
&= [-0.330, \, 1.716]
\end{align*}

\[
\boxed{CI_{95\%}(\beta_1) = [-0.330, \, 1.716]}
\]

\textbf{Step 6: 95\% CI for OR}

\begin{align*}
CI_{95\%}(OR) &= [\exp(-0.330), \, \exp(1.716)] \\
&= [0.719, \, 5.560]
\end{align*}

\[
\boxed{CI_{95\%}(OR) = [0.719, \, 5.560]}
\]

\textbf{Step 7: Hypothesis Test}

Test $H_0: \beta_1 = 0$ (equivalently, $OR = 1$):

\[
z = \frac{0.6931}{0.5218} = 1.328, \quad p\text{-value} = 2P(Z > 1.328) \approx 0.184
\]

\textbf{Conclusion:} Not statistically significant at $\alpha = 0.05$ (since $p = 0.184 > 0.05$ and CI includes 1).

\textbf{Summary Table:}

\begin{center}
\begin{tabular}{|l|c|c|c|c|}
\hline
\textbf{Parameter} & \textbf{Estimate} & \textbf{SE} & \textbf{95\% CI} & \textbf{p-value} \\
\hline
$\beta_0$ & 0.406 & -- & -- & -- \\
$\beta_1$ & 0.693 & 0.522 & [-0.330, 1.716] & 0.184 \\
OR & 2.000 & -- & [0.719, 5.560] & 0.184 \\
\hline
\end{tabular}
\end{center}
\end{solution}

\begin{explanation}
\textbf{Interpretation:}

Despite $OR = 2.0$ (doubling of odds), the result is not statistically significant due to:
\begin{itemize}
\item Small sample size ($n=70$)
\item Large SE (0.522) relative to estimate (0.693)
\item Wide CI including null value (OR = 1)
\end{itemize}

\textbf{Practical vs. Statistical Significance:}

The observed risks are $P(Y=1|X=0) = 18/30 = 0.60$ and $P(Y=1|X=1) = 30/40 = 0.75$.

Risk difference: $0.75 - 0.60 = 0.15$ (15 percentage points)

Relative risk: $RR = 0.75/0.60 = 1.25$

Note: $RR = 1.25 \neq OR = 2.0$ because baseline risk is high (60\%). OR and RR only agree when outcomes are rare.

\textbf{Power:} The smallest cell ($d=10$) contributes most to variance ($0.10$ of $0.272$), showing importance of adequate cell counts.
\end{explanation}

\newpage
